\documentclass[11pt,a4paper]{article}
\usepackage[a4paper,margin=27mm]{geometry}
\usepackage[T1]{fontenc}
\usepackage{lmodern,amsmath,amssymb,hyperref}
\newcommand{\E}{\mathbb{E}}
\newcommand{\dif}{\mathop{}\!\mathrm{d}}
\setlength{\emergencystretch}{2em}
\title{Release-v4 explanatory supplement (E04)}
\author{}
\date{18 October 2026}
\begin{document}\maketitle
\section{Compensated response and likelihood-product explanations}
\label{sec:vfour_math_explanations}
These finite-horizon explanations expand the retained manuscript identities;
they do not replace the scientific procedures or establish a stationary
nonlinear exact-memory theorem.

\subsection{Compensator cancellation}
With physical productivity $\E \kappa=n$, scalar initialization
$X_0\Phi$, zero initial excitation and no additional prehistory, set
$f=X_0\Phi+\rho(1*\Phi)$ and write $X$ for post-event stress in ordinary
time convolutions. Then $X=f-\Phi*\dif N_g$ and
$\lambda=\Psi(X_-)+h*\dif N_\kappa$.
The deterministic resolvent identities are
\[
Q_\gamma+nR_\gamma*h=R_\gamma,\qquad
Q_\gamma+\gamma R_\gamma*\Phi=\Phi.
\]
Substituting $\dif M_g=\dif N_g-m_1\lambda\,\dif t$
and $\dif M_\kappa=\dif N_\kappa-n\lambda\,\dif t$
into the right-hand side $S$ of the manuscript's compensated identity gives
\[
\begin{aligned}
S={}&X_0Q_\gamma+\rho(1*Q_\gamma)+\gamma R_\gamma*X
-m_1R_\gamma*\Psi-Q_\gamma*\dif N_g\\
&+m_1Q_\gamma*\lambda-m_1R_\gamma*h*\dif N_\kappa
+m_1nR_\gamma*h*\lambda\\
={}&X_0Q_\gamma+\rho(1*Q_\gamma)-Q_\gamma*\dif N_g
+\gamma R_\gamma*X.
\end{aligned}
\]
Here the first kernel relation combines the compensators, after which
$\lambda-\Psi=h*\dif N_\kappa$ cancels the remaining terms.
The second relation makes the first three final terms
$X-\gamma R_\gamma*X$, so $S=X$. This reasoning does not require
independence of the two mark martingales. Predictable intensities still
use the left limit; changes on the countable jump set do not affect
ordinary time convolutions. The algebra is causal on finite horizons,
without a transform assumption for the nonlinear sample path.

\subsection{An event-free deterministic check}
For $n=0$, $m_1=\ell=\gamma=\tau=X_0=\rho=1$ and
$\alpha=1/2$, take $\Psi(X)=X$ on the positive range of the path. Then
\[
\Phi(t)=e^t\operatorname{erfc}(\sqrt t),\qquad
X(t)=\Phi(t)+\int_0^t\Phi(u)\,\dif u,\qquad
\widehat Q(s)=\widehat R(s)=\frac1{s+\sqrt s+1}.
\]
Although no raw event occurs, $\dif M_g=-X\,\dif t$ and
$-R*\dif M_g=R*X$. The nonlinear residual is zero, giving
\[
X=Q+1*Q+R*X.
\]
At $t=1$, deterministic evaluation gives $X(1)=0.983546319407127$,
$Q(1)+(1*Q)(1)=0.623004105263315$, and
$(R*X)(1)=0.360542214143811$, up to displayed rounding.
The compensator accounts for the apparent discrepancy. Replacing only
the initial/loading kernels $Q$ by $\Phi$ instead adds $\gamma R*f$.
The linear link is used only on this positive stress range, and a positive
extension outside it can be supplied; no globally negative intensity is
needed. These are fixed-parameter evaluations, not simulated catalogues.

\subsection{The likelihood-product envelope}
At fixed generating nuisances, let
\[
\mathcal B=\int_0^TB(t)\,\dif t,\quad
J_\eta=\int_0^TB(t)(1-e^{-\eta G_s(t)})\,\dif t,\quad
q_i=\frac{B_i e^{-\eta G_i}+c_i}{B_i+c_i},\quad c_i=nH_i.
\]
At actual reference events $B_i+c_i>0$ almost surely. On
$\eta\in[\eta_-,\eta_+]\subset(0,\infty)$, the bound
$\sup_{z\ge0}ze^{-\eta_-z}=1/(e\eta_-)$ gives
$|q_i'|\le1/(e\eta_-)$ and
$|J_\eta'|\le\mathcal B/(e\eta_-)$.
As $0\le q_i\le1$ and $0\le J_\eta\le\mathcal B$,
\[
\begin{aligned}
|\partial_\eta L_\eta|
&\le e^{J_\eta}\left(|J_\eta'|\prod_iq_i+
\sum_i|q_i'|\prod_{j\ne i}q_j\right)\\
&\le\frac{e^{\mathcal B}(\mathcal B+N(T))}{e\eta_-}.
\end{aligned}
\]
No event factor is divided out. Since
$\E _0N(T)\le\mathcal B/(1-n)$, this is an integrable envelope
for differentiating the expectation of $L_\eta I$ with the fixed bounded
rejection indicator $I$. It is a derivative bound for $L_\eta$, not
$\log L_\eta$. It requires neither integrability of $\sum_iG_i$ nor
an added stress-mark moment assumption. The existing null-set convention,
fixed-rule conditions, positive-$\eta$ restriction and non-Wilks limitations
remain in force.

\subsection{Why the small-set clock is shared}
For one common Poisson measure with intensity
$\dif s\,\dif u\,F(\dif m)$, the no-proposal event in
$[0,t]\times[0,\Lambda_R]\times\mathsf M$ has probability
$e^{-\Lambda_Rt}$ because $F(\mathsf M)=1$. Both event-free flows from the
specified sublevel set have rates at most $\Lambda_R$. The first accepted
jump would therefore contradict this event. This justifies the retained
common-clock lower bound without asserting a rate bound after arbitrary
jumps.

\subsection{Paired MCSE convention}
For paired binary decisions let $D_i\in\{-1,0,1\}$ be first minus second,
with $b$ positive and $c$ negative discordances among $N_{\rm pair}$ pairs.
Since $\sum D_i=b-c$ and $\sum D_i^2=b+c$,
\[
\overline D=\frac{b-c}{N_{\rm pair}},\qquad
\widehat{\operatorname{MCSE}}=
\sqrt{\frac{b+c-N_{\rm pair}\overline D^2}{N_{\rm pair}(N_{\rm pair}-1)}}.
\]
This is sample standard deviation divided by $\sqrt{N_{\rm pair}}$.
At 150 pairs, 29/21 yields 4.709613821 percentage points and 61/6 yields
4.577609808. The empirical-variance divisor $N_{\rm pair}$ instead of
$N_{\rm pair}-1$ yields 4.693888856 and 4.562325593 from the same outcomes.
These are conditional descriptive summaries under fixed-bank and adaptive
reuse qualifications, without a new test or confidence interval.

\end{document}
